

\documentclass{article}
\usepackage[utf8]{inputenc}
% Better font encoding for French and scalable fonts
\usepackage[T1]{fontenc}
\usepackage{lmodern}
\usepackage{amsmath, amssymb, amsfonts, amsthm, stmaryrd}
\usepackage{tikz}
\usetikzlibrary{matrix,arrows}
\usepackage{tikz-cd}
\usepackage{geometry}
\usepackage[french]{babel}
\usepackage{amsopn}

\title{Une méthode solide de test statistique pour détecter les manoeuvres}
\author{Etienne Chevrolat}
\date{}

\begin{document}

\maketitle

\begin{abstract}
    On cherche ici à expliquer une méthode statistique avancée pour tester la présence de ruptures au sein des séries de paramètres orbitaux.
    On teste à la fois des ruptures de niveau (manoeuvres impulsionnelles), et des ruptures de pentes (manoeuvres type propulsion électrique).
\end{abstract}

\section{Le modèle statistique}

\subsection{Le cadre}

On va considérer $y_{sat}$ un paramètre orbital issu d'un TLE (demi-grand axe ou inclinaison par exemple) sur un historique long (2000 TLEs par exemple).

On découpe cet historique long en fenêtres courtes de $N = 20$ TLEs typiquement. 

Considéront donc une de ces fenêtres. On réindexe le temps sur le milieu de la fenêtre : 
$ 
\forall k \in [1, \dots, N] : \tau_k = t_k - t_{centre}
$
L'instant central $t_{centre}$ est l'instant où a lieu la manoeuvre s'il y en a une dans cette fentêtre temporelle.

Sur cette fenêtre temporelle courte, on va approximer $y_{sat}$ par son polynôme de Taylor à l'ordre $p$ ainsi que par des termes de ruptures: 

$$
y_{sat}(\tau) = c_0 + c_1 \tau ^1 + \dots + c_p \tau ^ p + A_{niveau} \mathbb{1}_{\tau > 0} + B_{pente} max(0, \tau) + \epsilon 
$$

\vspace{0.3cm}
\noindent\fbox{
\begin{minipage}{0.98\textwidth}
\textbf{Modèle linéaire gaussien}
On considère un bruit gaussien : $\epsilon \sim \mathcal{N}(0,I)$
On peut réécrire l'approximation précédente sous la forme :
$
y_sat = X \beta + \epsilon
$
\begin{itemize}
    \item 
    X = \begin{pmatrix}
        a & b \\ c & d 
        \end{pmatrix}
\end{itemize}
\end{minipage}
}




On peut ensuite lier des variations de vitesse dans le référentiel inertiel terrestre aux variations des paramètres orbitaux de l'ellipse, et vice-versa. 

\vspace{0.3cm}

\noindent\fbox{
\begin{minipage}{0.98\textwidth}
\textbf{Variations de vitesse et variations de paramètres orbitaux}\newline
On considère une variation $\Delta v = (\Delta v _r , \Delta v _u ,\Delta v _h )$ de la vitesse décomposée dans le référentiel terrestre en vitesses radiale, tangentielle et normale.
Les équations de Gauss nous donnent dans la limite quasi-circulaire:  
\begin{itemize}
    \item $\Delta a = \dfrac{2}{n} \Delta v _u$

\end{itemize}
\end{minipage}

}
\vspace{0.3cm}

On en déduit les interprétations suivantes : 
Des variations de a et e traduisent des manoeuvres \textit{in plane}  (variation de vitesse radiale)
Des variations de i et $\Omega$ traduisent des  manoeuvres \textit{out plane} (variation de vitesse normale)

\subsection{Two Line Elements}

Les paramètres orbitaux décrits ci-dessus ne sont pas mesurés directement : ils sont diffusés par les organismes de surveillance de l'espace sous la forme de \textit{Two Line Elements} (TLE).
Un TLE est un format standardisé qui encode, sur deux lignes de 69 caractères, l'ensemble des éléments nécessaires pour reconstituer l'orbite d'un objet à un instant de référence appelé \textit{epoch}.

\vspace{0.3cm}
\noindent\fbox{%
\begin{minipage}{0.98\textwidth}

\textbf{Contenu d'un TLE}
\begin{itemize}
    \item \textbf{Identification} : numéro de catalogue NORAD, classification, désignateur international (année de lancement, numéro de lancement) et numéro du jeu d'éléments.
\end{itemize}
\end{minipage}

}
\vspace{0.5cm}


\end{document}
 